\subsection{Внешняя и внутренняя меры Жордана и Лебега}
\begin{note}
	Всюду в этом параграфе любое рассматриваемое множество будет являться подмножеством $K_I=\l[-\frac12,\frac12\r]^n$.
\end{note}
\begin{definition}
	\textit{Внешней мерой Жордана} множества $A$ называется: \[\mj(A):=\inf_{A\subset\bigcup_{i=1}^mM_i}\sum_{i=1}^m|M_i|\]
\end{definition}
\begin{definition}
	\textit{Внешней мерой Лебега} множества $A$ называется:
	\[\m(A):=\inf_{A\subset\bigcup_{i=1}^\i M_i}\sum_{i=1}^\i|M_i|\]
\end{definition}
\begin{note}
	Запись $\ma$ означает, что выражение справедливо и для меры Жордана, и для меры Лебега.
\end{note}
\begin{theorem}\nf{(Основные свойства внешних мер)}\label{20_th1}
	\begin{enumerate}
		\item Для любого элементарного множества $M$:
		\[\mj(M)=\m(M)=|M|\]
		\item Если $\ds A\subset\bc_{i=1}^\i A_i$, то:
		\[\ma(A)\le\sum_{i=1}^r\ma(A_i)\]
		\item Для любого множества $A\subset K_I$:
		\[\m(A)\le\mj(A)\]
		\item Для любого $\ds A\subset\bc_{i=1}^\i A_i$:
		\[\m(A)\le\sum_{i=1}^\i\m(A_i)\]
	\end{enumerate}
\end{theorem}
\begin{proof}
	\begin{enumerate}
		\item \[M\subset M\Ra\ma(M)\le|M|\]
		\[M\subset\bc_iM_i\Ra|M|\le\sum_i|M_i|\Ra|M|\le\ma(M)\]
		\item \[A_i\subset\bc_jM_{i,j}\Ra A\subset\bc_{i=1}^r\bc_jM_{i,j}=1\Ra\ma(A)\le\sum_{i=1}^r\sum_j|M_{i,j}|\le\sum_{i=1}^r\ma(A_i)\]
		\[\sum_j|M_{i,j}|\le\ma(A_i)+\frac\eps2\]
		\item Следует из пункта 2.
		\item \[A_i\subset\bc_{j=1}^\i M_{i,j}\Ra A\subset\bc_{i=1}^\i\bc_{j=1}^\i M_{i,j},\ \sum_{j=1}^\i|M_{i,j}|\le\m(A_i)+\frac\eps2\]
	\end{enumerate}
	Получили требуемое.
\end{proof}
\begin{definition}
	Дополнением множества $A$ назовём $A':=K_I\sm A$.
\end{definition}
\begin{definition}
	Внутренней мерой Лебега (Жордана) множества $A\subset K_I$ называется:
	\[\mu_*^{(J)}(A):=1-\ma(A')\]
\end{definition}
\begin{theorem}\nf{(Основные свойство внутренней меры)}\label{20_th2}
	\[\fa A\subset K_I\ \mu_*^{(J)}(A)\le\mu^*_{(J)}(A)\]
\end{theorem}
\begin{proof}
	\[A\subset\bc_iM_i,\ A'\subset\bc_jN_j\Ra A\sqcup A'=K_I\subset\l(\bc_iM_i\r)\cup\l(\bc_jN_j\r)\]
	\[|K_I|=1\le\sum_i|M_i|+\sum_j|N_j|\Ra1\le\ma(A)+\ma(A')\]
\end{proof}
\begin{example}
	Пусть $A=\Q\cap K_I$ при $n=1$. \\
	Для меры Жордана можем покрыть весь отрезок брусьями, а для меры Лебега в качестве элементарных множеств возьмём точки $\in\Q$. Их счётно, объём каждой 0. Тогда имеем:
	\[\mj(A)=1,\ \m(A)=0,\ \mu_*^J(A)=1,\ \mu^*(A)=0\]
\end{example}
\subsection{Измеримое множество. Основные свойства мер Жордана и Лебега}
\begin{note}
	В этом параграфе всё ещё работаем с $K_I$.
\end{note}
\begin{definition}
	Множество $A$ называется \textit{измеримым по Лебегу (Жордану)}, если $\ma(A)=\mu_*^{(J)}(A)$, их общее значение называется мерой Лебега (Жордана) множества $A$.
\end{definition}
\begin{note}
	Из (\ref{20_th1}) следует, что элементарные множества измеримы, при этом:
	\[\mu_{(J)}(M)=|M|\]
\end{note}
\begin{theorem}\nf{(Критерий измеримости)}\label{20_th3}
	Множество $A\subset K_I$ измеримо по Лебегу (Жордану) тогда и только тогда, когда:
	\[(\fa\eps>0)(\ex\text{ элементарное } M_\eps)\ \mu^*_{(J)}(A\tr M_\eps)<\eps\]
\end{theorem}
\begin{proof}
	Достаточность:
	\[K_I=A\sqcup A',\ \ma(K_I)=|K_I|=1\le\ma(A)+\ma(A')\]
	\[K_I=M_\eps\sqcup M_\eps'\Ra|K_I|=1=|M_\eps|+|M_\eps'|\]
	\[A\subset M_\eps\sqcup(A\sm M_\eps)\Ra\ma(A)\le\ma(M_\eps)+\ma(A\tr M_\eps)<|M_\eps|+\eps\]
	\[A'\subset M_\eps'\sqcup(A'\sm M_\eps')=M_\eps'\sqcup(M_\eps\sm A)\Ra\ma(A')\le|M_\eps'|+\eps\]
	Необходимость:
	$A$ измеримо, тогда имеем:
	\[\ma(A)+\ma(A')=1\]
	\[(\fa\eps>0)\ A\subset\bsc_iP_i,\ A'\subset\bsc_iQ_i,\ \sum_i|P_i|<\ma(A)+\eps,\ \sum_j|Q_j|<\ma(A')+\eps\]
	Для меры Лебега:
	\[(\ex r\in\N)\ \sum_{i=r+1}^\i|P_i|<\eps,\ \sum_{j=r+1}^\i|Q_j|<\eps\]
	Положим:
	\[P_\eps:=\bsc_i^{(r)}P_i,\ Q_\eps:=\bsc_j^{(r)}Q_j,\ A\tr P_\eps=(A\sm P_\eps)\cup(P_\eps\sm A)\]
	Хотим доказать:
	\[\ma(A\tr P_\eps)<6\eps\]
	\[A\sm P_\eps\subset\bc_{i=r+1}^\i P_i\Ra\ma(A\sm P_\eps)\le\sum_{i=r+1}^\i|P_\eps|<\eps,\text{ для меры жордана }\emptyset\]
	\[P_\eps\sm A\subset(P_\eps\cap Q_\eps)\cup(Q_\eps'\sm A)=(A'\sm Q_\eps)\cup(P_\eps\sm Q_\eps)\]
	\[\ma(A'\sm Q_\eps)<\eps,\text{ для меры Жордана }\ma(A'\sm Q_\eps)=0\]
	\[|P_\eps\cap Q_\eps|=|P_\eps|+|Q_\eps|-|P_\eps\cup Q_\eps|\]
	\[|P_\eps|\le\sum_i|P_i|<\ma(A)+\eps,\ |Q_\eps|\le\sum_j|Q_j|<\ma(A)+\eps\]
	\[K_I\subset\l(\bsc_iP_i\r)\cup\l(\bsc_jQ_j\r)=P_\eps\cup\l(\bsc_{i=r+1}^\i P_i\r)\cup Q_\eps\cup\l(\bsc_{j=r+1}^\i Q_j\r)\]
	\[1=|K_I|\le|P_\eps\cup Q_\eps|+\sum_{i=r+1}^\i|P_i|+\sum_{j=r+1}^\i|Q_j|<|P_\eps\cup Q_\eps|+2\eps\Ra|P_\eps\cup Q_\eps|>1-2\eps\]
	Имеем:
	\[|P_\eps\cap Q_\eps|<\ma(A)+\eps+\ma(A')+\eps-1+2\eps=4\eps\]
	Получили требуемое.
\end{proof}
\begin{theorem}\label{20_th4}
	Семейство измеримых по Лебегу (Жордану) подмножеств $K_I$ образует алгебру множеств.
\end{theorem}
\begin{proof}
	\[(\fa\eps>0)\ \mu_{(J)}(A\tr P_\eps)<\eps,\ \ma(B\tr Q_\eps)<\eps\]
	\begin{multline*}
		(A\cup B)\tr(P_\eps\cup Q_\eps)=((A\cup B)\sm (P_\eps\cup Q_\eps))\cup((P_\eps\cup Q_\eps)\sm(A\cup B))\subset\\
		\subset(A\sm P_\eps)\cup(B\sm Q_\eps)\cup(P_\eps\sm A)\cup(Q_\eps\sm B)=(A\tr P_\eps)\cup(B\tr Q_\eps)\Ra\m((A\cup B)\tr(P_\eps\cup Q_\eps))\le\m(A\tr P_\eps)+\m(B\tr Q_\eps)<2\eps
	\end{multline*}
	Тогда по (\ref{20_th3}) $A\cup B$ измеримо.
	\[(A\cap B)'=A'\cup B',\ A\text{ - измеримо }\lra A'\text{ - измеримо}\]
	\[A\sm B=A\cap B',\ A\tr B=(A\cup B)\sm(A\cap B)\]
	Получили требуемое.
\end{proof}